% Copyright 2026 icekylinx and GPT-6 Astra Pro. Apache-2.0.
% Curated from sat_research_checkpoint.zip: paper/direct_mixed_crossing_transfer.tex
% Complete simultaneous-insertion and ordered-crossing lemmas retained; downstream field/older numerical branches omitted.
\subsection{A direct noncommutative transfer in every characteristic}
\label{sec:direct-mixed-crossing}

The full crossing identity contains more information than its
commutative sum-of-squares contraction. Keeping that identity permits a
direct transfer with no division by two. In particular, the integral
Schur construction also gives noncommutative circuit bounds in
characteristic two.

\begin{lemma}[Simultaneous insertion without division]
\label{lem:simultaneous-insertion}
Let $R$ be a commutative unital ring. Suppose a division-free
commutative circuit of size $C$ simultaneously computes bilinear forms
$F_i(X,Y)$, and a division-free noncommutative circuit of size $D$
computes polynomials $G_i$. If
$F_i=\sum_{a,b}f_{i,a,b}X_aY_b$, then
\[
 \sum_i\sum_{a,b}f_{i,a,b}X_aG_iY_b
\]
has a division-free noncommutative circuit of size $O(C+D+t)$, where
$t$ is the number of forms. The statement is valid over $R$, including
in characteristic two and in the presence of zero divisors.
\end{lemma}

\begin{proof}
Keep the four bidegree components $(0,0),(1,0),(0,1),(1,1)$ of the
first circuit; this costs a constant factor and uses no division.
Its bilinear multiplication gates can be written
$A_\ell(X)B_\ell(Y)$, where both factors are linear. Regard each such
product as a new input $W_\ell$. The remaining bilinear computation is
a scalar-linear circuit
\[
 F_i=\sum_\ell c_{i,\ell}W_\ell.
\]
Transpose only this scalar-linear circuit to compute
$H_\ell=\sum_i c_{i,\ell}Z_i$. Explicitly, initialize the adjoints of
the output gates by $Z_i$, propagate an adjoint to both inputs of an
addition gate, and multiply by the same scalar when traversing a scalar
multiplication gate backwards. Sum the contributions arriving at each
gate. The number of operations is linear in the original number of
edges and the output count, hence $O(C+t)$.

Substitute $Z_i=G_i$ and output
$\sum_\ell A_\ell(X)H_\ell B_\ell(Y)$. Since every $c_{i,\ell}$
is central, this is the desired ordered insertion. The linear forms
and the transposed circuit remain shared. All identities use only
addition and multiplication, so the argument is characteristic-free.
\end{proof}

\begin{lemma}[Mixed crossing to an ordered identity polynomial]
\label{lem:direct-mixed-crossing}
Suppose arrays over $R$ satisfy
\begin{equation}
\label{eq:direct-mixed-crossing}
 [x=v][y=u]
 =\sum_{j\in J} f_{j,x,y}g_{j,v,u}
  +\sum_{j\in K} h_{j,x,u}k_{j,v,y}.
\end{equation}
If all associated bilinear forms are simultaneously computable by
division-free commutative circuits of total size $C$, then
\[
 \operatorname{ID}_{p,q}(X,Y)
   =\sum_{x\in[p],\,y\in[q]}X_xY_yX_xY_y
\]
has a division-free noncommutative circuit of size
$O(C+|J|+|K|)$.
\end{lemma}

\begin{proof}
Contract~\eqref{eq:direct-mixed-crossing} against the ordered word
$X_xY_yX_vY_u$. A row-oriented term is the product
\[
 \left(\sum_{x,y}f_{j,x,y}X_xY_y\right)
 \left(\sum_{v,u}g_{j,v,u}X_vY_u\right).
\]
A column-oriented term is
\[
 \sum_{x,u}h_{j,x,u}X_x
   \left(\sum_{v,y}k_{j,v,y}Y_yX_v\right)Y_u.
\]
Bihomogenization turns a commutative bilinear circuit into either
ordered bilinear circuit with constant overhead. When one shared
factor circuit has both row-oriented $XY$ outputs and column-oriented
$YX$ outputs, make two ordered copies and select the required outputs
from each. This costs a constant factor; no shared product is flipped
independently for different outputs. All column insertions
together are handled by Lemma~\ref{lem:simultaneous-insertion}; row
products and their sums add only the number of terms. The right side
therefore contracts to $\operatorname{ID}_{p,q}$ with the stated cost.
No polarization, square root of $-1$, or inversion of two is used.
\end{proof}
